\documentclass[11pt,a4paper]{article}
\usepackage[utf8]{inputenc}
\usepackage[T1]{fontenc}
\usepackage[french]{babel}
\usepackage{lmodern}
\usepackage[margin=2.2cm]{geometry}
\usepackage{longtable,booktabs,array}
\usepackage{xcolor}
\usepackage[hidelinks]{hyperref}
\usepackage[numbers,sort]{natbib}
\usepackage{textcomp}

\newcommand{\commit}[1]{\texttt{\small #1}}
\newcommand{\ecart}[1]{\textcolor{red!60!black}{#1}}
\setlength{\parskip}{0.4em}
\setlength{\parindent}{0pt}

\title{Journal de développement de \emph{vetsim} et de l'application \emph{VetTwins}\\[0.3em]
\large Modifications, vérifications, écarts et références}
\author{Projet VetTwin}
\date{Dernière mise à jour : 17 novembre 2026}

\begin{document}
\maketitle

\begin{center}
\fbox{\parbox{0.92\linewidth}{\textbf{Usage pédagogique uniquement.} Le moteur est un patient virtuel canin destiné à l'enseignement. Aucun résultat ne doit servir à décider d'une dose ou d'un traitement réel. Les paramètres sont partiellement calés sur des moyennes publiées et ne sont pas validés cliniquement.}}
\end{center}

\section{Règles du projet}
\begin{itemize}
  \item L'implémentation Python de \texttt{reference/python/} est \emph{normative}. Le moteur Rust doit la reproduire dans les tolérances du golden (\texttt{tests/golden/cases.json}, 962 valeurs).
  \item Ni les modules JSON, ni les cibles de calibration, ni le golden, ni les valeurs de paramètres ne sont modifiés pour faire passer un test. Tout écart est rapporté avec ses chiffres.
  \item Aucune donnée clinique ou pharmacologique n'est inventée. Une valeur manquante devient un placeholder marqué \texttt{evidence: "D"}.
  \item Les formules des modules passent uniquement par le DSL (arbre d'opérations fermé), jamais par \texttt{eval}.
  \item Déterminisme : graine explicite, générateur ChaCha8, fonctions mathématiques \texttt{libm} (mêmes bits sur toutes les plateformes).
  \item Les références bibliographiques sont vérifiées (OpenAlex, DOI ou PubMed) avant d'être citées. La section~\ref{sec:refs} donne le statut de chacune.
\end{itemize}

\section{Chronologie des modifications}
Chaque entrée indique le commit, ce qui a été fait, la vérification qui fait foi et les écarts constatés. Branche \texttt{m0-m1-workspace-dsl}.

\subsection*{17 novembre 2026}
\begin{longtable}{@{}p{1.5cm}p{4.2cm}p{10.3cm}@{}}
\toprule
\textbf{Commit} & \textbf{Étape} & \textbf{Contenu, vérification, écarts}\\
\midrule\endhead
\commit{d46cbe8} & Kit d'origine & Kit livré, inchangé. Référence Python : 962/962 valeurs golden conformes.\\
\commit{84ba719} & M0 & Workspace Cargo, lints, CI. \texttt{fmt}, \texttt{clippy -D warnings} et tests verts.\\
\commit{2454c63}, \commit{a9395fa} & M1 : DSL, modules & DSL fermé, validation des modules avec chemins JSON, patient sain, effets maladie. Égalité avec Python à $10^{-9}$ (5 poids, 4 maladies × 9 sévérités, 45 combinaisons, 44 expressions). CLI : 4161/4161 paramètres.\\
\commit{0e50c38} & M2 : gaz, acide-base & Stewart (albumine canine \citep{constable2005}), Hill et Bohr \citep{cambier2004}, rein, lactate. Golden 390/390 ; test strict 3250/3250 à $10^{-9}$.\\
\commit{2723d75}, \commit{6fd42bf} & M3 : hémodynamique & Dormand--Prince 5(4) \citep{dormand1980} à la place de LSODA, boucle baroréflexe. Golden 630/630, écart maximal $0{,}0075\times$ la tolérance. \ecart{Écart : MMVD $s=1{,}0$ (hors golden), critère d'arrêt du baroréflexe sur le fil (marge $6{,}7\times10^{-5}$), \texttt{lv\_edp} à $1{,}02\times$ la tolérance.} Décision : garder la référence telle quelle.\\
\commit{8356400} & M4 : PK/PD & PK à 2 compartiments et dépôt oral, effets JSON, diurèse \citep{koh2021,her2020}. Golden 962/962.\\
\commit{d9de070} & M5 : population & Graines explicites, banc de calibration. 160/160 lignes conformes, écart moyen 1,58\,\%. Critère adopté : écart relatif $<5$\,\% ou $<3$ erreurs-types combinées.\\
\commit{e2a5842} & Déterminisme & Noyau sur \texttt{libm}. Le cas MMVD $s=1{,}0$ bascule de l'autre côté du seuil : ce n'est pas une correction.\\
\commit{b0c0325} & M6 : autonome, ECG, PCG & VFC dans $\pm1$ ET des cibles \citep{bogucki2015} sur 20 graines. Rejeu des tirages NumPy : instants R à $3\times10^{-14}$\,s près.\\
\commit{b5ca64d}--\commit{de04c4b} & M7 : FFI, VetTwins & UniFFI, démo Kotlin Multiplatform. Sorties identiques au bit près entre macOS natif, la JVM via FFI et wasm32 (32/32).\\
\commit{8f0d0fa} & Respiration, schéma & Courbes respiratoires (volume = profil normalisé $\times$ volume courant alvéolaire VA/FR, \ecart{espace mort non modélisé}). Schéma des blocs du modèle et rattachement paramètre $\to$ bloc.\\
\commit{832437a} & Pressions & PAS et PAD instantanées = PAM instantanée $\pm$ écart du battement simulé. \ecart{Pression pulsée supposée constante.}\\
\commit{b7cccee}--\commit{790bdab} & Auscultation (a--d) & Port fidèle de \texttt{signals\_pcg\_v5.py} (signal à $4\times10^{-15}$ près), configuration JSON avec provenance, critères mesurés par \texttt{pcg\_metrics.py}, FFI, onglet Auscultation (son, spectrogramme, ECG et pressions alignés, zoom, curseur). Calage sur deux enregistrements \citep{tiho}. \ecart{Écarts de la référence conservés : B1$\to$B2 0,275\,s à FC 90, B2 $>$ B1 dans 75\,\% des paires, diastole/systole 1,24 et 1,60.} 100 graines au lieu de 20 (appariement B1/B2 instable).\\
\commit{9d3f7b7}, \commit{e59a51c} & PK/PD répétée & Examen à un instant quelconque d'un schéma ; effets E(t) ; indicateurs au choix sur 1 à 14 jours. \ecart{Furosémide 2\,mg/kg per os toutes les 12\,h : K$^+$ de 4,1 à environ 2,4--3,0\,mmol/L (PD placeholder du module).}\\
\commit{8ffcfac} & Journal & Création de ce journal et de la bibliographie vérifiée.\\
\commit{12e0e4a} & Édition des doses & Modification d'un schéma existant ; ajustement prise par prise (dose modifiée, prise omise = 0). Vérifié : prise omise à 48\,h $\Rightarrow$ concentration $<5$\,\% du pic habituel ; prise doublée à 72\,h $\Rightarrow$ pic $>1{,}8\times$.\\
\commit{1fe10af} & Application web & API JSON déplacée dans \texttt{vetsim-api} (partagée), liaison WebAssembly \texttt{vetsim-wasm} (\texttt{wasm-bindgen} 0.2.100), interface web statique (5 onglets, cas partageables par l'adresse). Vérifié : 882/882 valeurs identiques au bit près entre Chrome (WebAssembly) et le binaire natif ; captures des onglets relues.\\
\commit{3d8d5c8} & Vérification bibliographique & Comparaison chiffrée du modèle à 30 sources (section~\ref{sec:verif}, \texttt{docs/VERIFICATION\_LITTERATURE.md}). Aucun module ni paramètre modifié.\\
\commit{db2ac1e} & Proposition toxicité & \texttt{docs/PROPOSITION\_TOXICITE.md} : termes toxiques séparés $T$ dans la PD (au lieu de $E>1$, qui ne fait qu'amplifier l'effet thérapeutique), critères de décompensation sourcés, risque de décès par risque instantané (flux ChaCha distinct). En attente de validation, rien d'implémenté.\\
\commit{266add1} & Toxicité et décompensation & Validé par l'utilisateur : termes toxiques séparés, placeholders D, référence Python étendue, \og tout mécanistique\fg{}. PD : termes toxiques $T$ (opioïde : apnée ; bêtabloquant : dépression myocardique et sinusale ; pimobendane : vasoplégie). L'effet thérapeutique reste plafonné (occupation des récepteurs) : c'est la réponse retenue à la demande \og désactiver la maximisation\fg{}. Point fixe hémodynamique--métabolisme : acidémie, hypoxie tissulaire, ischémie coronaire $\to$ contractilité ; acidémie $\to$ vasoplégie ; hypokaliémie $\to$ faiblesse ventilatoire. Décès = absence d'état viable (arrêt cardiaque, effondrement circulatoire, VO\textsubscript{2}~$>$~DO\textsubscript{2}, pas d'état stationnaire, spirale). 11 paramètres placeholders D. Vérifié : golden 962/962 (Python et Rust), paramètres 4788/4788, calibration 164/164 (aucun couplage actif sur 264 cas, écart identique à la référence d'origine), fixture Rust~=~Python, 3 mutations sur 3 détectées. \ecart{Cas golden avec médicament : écart de 0,054\,\% au plus (queue toxique à la CE50). MMVD sévère létale chez le chien de 3\,kg (l'ancien modèle donnait déjà ER~1,56 et PvO\textsubscript{2} au plancher) ; un chien de 60\,kg reste \og stable\fg{} en MMVD $s=1$.}\\
\commit{69e62b2} & Interfaces & Concentrations fixes retirées (web, bureau, Android) : tout passe par la PK/PD. Bandeau d'issue, vues \og patient décédé\fg{}, bandes d'état sur les courbes du traitement. Décès absorbant dans \texttt{timecourse} et à l'examen. Auscultation : arrêter, rejouer, boucle, fenêtre seule. Vérifié : 882/882 valeurs identiques au bit près (navigateur), 8 tests KMP dont un surdosage mortel, APK construit.\\
\commit{87ca549} & Proposition rythme & \texttt{docs/PROPOSITION\_RYTHME.md} : troubles du rythme modélisés mécaniquement (conduction par événements : nœud sinusal, nœud AV à délai dépendant de la récupération et période réfractaire, échappements, foyer ectopique, FA), hémodynamique battement par battement avec oreillettes contractiles, ECG et sons dérivés. Ce sont de nouvelles maladies (décision de l'utilisateur). Recherches bibliographiques en cours (rythme, sons cardiaques réels). Rien d'implémenté.\\
\commit{965df07} & Oreillettes (structure) & Élastance active sur \texttt{pv} (OG) et \texttt{sv} (OD), déclenchée PR avant le ventricule ; mesure des ondes E et A mitrales. Gain nul en attendant le calage : résultats identiques (golden 962/962, paramètres 5016/5016). \ecart{Calage impossible dans les tolérances du golden : dès 0,8\,\% de contribution atriale, 5 valeurs sortent (PTDVG MMVD $s=0{,}8$). Décision de l'utilisateur : nouvelle référence, calée sur le rapport E/A du chien sain (recherche en cours).}\\
\commit{bd51828} & Couche de conduction & \texttt{rhythm.py} (référence Python) : nœud sinusal, oreillettes (FA, arrêt atrial hyperkaliémique avec conduction sino-ventriculaire), nœud AV (période réfractaire, délai de récupération, conduction cachée), His-Purkinje, échappements, foyer ectopique. Calés : FA 129,6/min et CVRR 0,291 (cibles 129,7 et 0,309, $n=50$) ; échappement 42/min ($n=97$) ; arrêt atrial à K~8,5 (Hariman 1983, résumé). 11 critères avec contrôles d'échec. Pas encore branchée au moteur. Hyperkaliémie (remarque de l'utilisateur) : bradycardie seulement au-delà de K~$\approx$~10, cohérent avec 9--21\,\% de chiens bradycardes en crise addisonienne (Chalifoux 2023, $n=84$).\\
\commit{3e83f20} & Volume de l'auscultation & Réglage $\times1$ à $\times8$ (web : gain et compresseur ; KMP : gain et limiteur doux tanh).\\
\commit{a29c017} & Conduction en Rust & \texttt{rhythm.rs} : 13 scénarios identiques à la référence à $10^{-12}$ (tirages rejoués), 3 mutations sur 3 détectées.\\
\commit{1a6255e}, \commit{393d4e8} & Hémodynamique battement par battement & \texttt{simulate\_hemo\_beats} (Python, puis Rust) : activation propre à chaque battement (systole $0{,}37\sqrt{RR}$), contraction atriale à chaque activation organisée, indicateurs par battement. Séquence régulière = simulation périodique à 0,6\,\% près. Avec l'oreillette calée sur E/A ($k_{la}=1{,}45$), la FA à 131/min abaisse le débit de 17\,\% par rapport au sinusal à la même FC (\textbf{$-18$\,\% rapporté} chez le chien, Naito 1983, résumé) ; le déficit de pouls émerge (VES de 12,6 à 50\,mL). Rust~=~Python à 1\,\% (moyennes) et 2\,\% (VES par battement), 3 mutations sur 3 détectées. Pas encore branchée au moteur.\\
\commit{(ce commit)} & Nouvelle référence & Décision de l'utilisateur (l'oreillette ne tenait pas dans l'ancien golden). Conduction branchée au moteur (sinusal 1:1 : périodique avec le PR de la conduction ; sinon battement par battement, aléa SplitMix64 identique Python/Rust) ; oreillette active ($k_{la}=1{,}45$, E/A) ; \texttt{kfac} supprimé (K$^+$ par la conduction) ; FC rapportée = FC ventriculaire ; ECG et battements du son construits depuis la conduction (P selon l'activation atriale, PR par battement, QRS large et T discordante, QT raccourci en hyperkaliémie) ; échelle unique de l'ECG. Golden et calibration régénérés : Rust 962/962, calibration 164/164, cohorte saine dans toutes les cibles ; décès MMVD stade D 2/24 (Python), 3/96 (Rust). \ecart{MMVD stade C moins congestif (POG s 0,8 : 20,6 $\to$ 17,4) ; Addison sans frein : FC 141 à K 6,8 ; PAM saine 106 (cible 103).} Détail : \texttt{docs/NOUVELLE\_REFERENCE.md}.\\
\bottomrule
\end{longtable}

\section{Décisions et écarts connus}
\begin{enumerate}
  \item \textbf{Critère d'arrêt du baroréflexe.} La sortie de la référence est discontinue près du seuil $|d_t-d|<0{,}04$. La référence est conservée (décision du 17/11/2026), voir \texttt{docs/SOLVER.md}.
  \item \textbf{Durée de systole trop longue.} Avec $t_s = 0{,}37\sqrt{T}$, B1$\to$B2 vaut 0,27\,s à FC 90 et les rapports diastole/systole sont trop faibles. La cause est hémodynamique, pas audio.
  \item \textbf{Hypokaliémie sous furosémide répété.} Elle découle de la PD du module (\texttt{k} $-= 0{,}04\,U$), qui est un placeholder. À examiner lors de la vérification bibliographique.
  \item \textbf{Définition de la diurèse du diurétique} (\texttt{pkpd.py}). Débit induit $u(t)=0{,}25\,E(C)$\,mL/kg/min (supplément dû au médicament ; la diurèse de base n'est pas modélisée). \og Diurèse effective cumulée\fg{} $U(t)=\int 60\,u(s)\,e^{-(t-s)/24\,\mathrm{h}}\,ds$ (mL/kg) : volume induit cumulé avec oubli exponentiel. Ce n'est pas le volume d'urine émis ; il fixe les pertes d'électrolytes (K$^+ - 0{,}04\,U$, Cl$^- - 0{,}5\,U$, Na$^+ - 0{,}3\,U$). Déficit $D$ : $\dot D = 0{,}3\,u\,\mathrm{bw}\,60 - D/4\,\mathrm{h}$ (mL). Placeholders.
  \item \textbf{Arythmie sinusale respiratoire à fréquence élevée} (remarque de l'utilisateur, 17/11/2026 : la RSA est attendue au repos et à basse fréquence). Mesure (5 graines, 120\,s) : sain FC 88, RSA 46\,bpm ; MMVD C FC 122, RSA 38\,bpm ; MMVD D FC 161, RSA 22\,bpm ; atropine 0,01\,mg/L FC 142, RSA 17\,bpm. La RSA diminue avec le tonus vagal mais reste trop ample en tachycardie. Causes (\texttt{autonomic.py}) : modulation respiratoire proportionnelle au tonus vagal moyen, effet relatif sur la FC (même variation relative, davantage de bpm à FC élevée). Correction possible (non faite, modèle normatif) : modulation non linéaire du tonus vagal ou effet vagal absolu sur la FC, à fonder sur la littérature (perte de l'arythmie sinusale dans l'insuffisance cardiaque canine).
  \item \textbf{Approximations d'affichage.} Pressions alignées sur le son : le battement simulé est posé à chaque battement. Volume respiratoire : sans espace mort. PAS et PAD instantanées : pression pulsée constante.
  \item \textbf{Placeholders (evidence D).} PD de tous les médicaments ; PK de l'atropine, du bêtabloquant et de l'opioïde ; module pneumonie ; site d'auscultation \og apex\fg{} ; seuils de grade des souffles.
\end{enumerate}

\section{Vérification bibliographique (17 novembre 2026)}\label{sec:verif}
Cinq recherches, puis une sixième sur l'ionogramme de la MRC, ont été menées par des sessions IA séparées de ce projet. Ce n'est pas une relecture externe. Les valeurs du modèle viennent de \texttt{vetsim-cli} en mode déterministe (20\,kg). Accès : \emph{lu} (texte de la source), \emph{résumé}, \emph{moteur} (résumé de moteur de recherche). Le détail figure dans \texttt{docs/VERIFICATION\_LITTERATURE.md}.

\textbf{Conformes.}
\begin{itemize}
  \item Chien sain \citep{haskins2005,cambier2004,constable2005} (résumés) : FC 87,6 contre 87 ; PAM 101 contre 103 ; PAP 13,9 contre 14 ; pH 7,38 ; PaCO\textsubscript{2} 40,2 ; P50 30 ; Bohr $-0{,}37$ ; A\textsubscript{tot} 0,469\,mmol/g.
  \item VFC du chien sain \citep{bogucki2015} (lu) : SDNN 185 contre $208{,}9\pm77{,}1$\,ms ; LF/HF 0,23 contre $0{,}28\pm0{,}11$.
  \item Fraction de régurgitation en B1 et C \citep{kawai2026}.
  \item FC et PAS en ICC \citep{yin2026}.
  \item Pimobendane per os \citep{her2020} : C\textsubscript{max} 43,8 contre $49{,}1\pm28{,}7$\,ng/mL ; ASC 149 contre $148\pm72$.
  \item Addison \citep{adler2007,baumstark2014,hall2023} : Na:K, Na$^+$ et K$^+$ en crise.
  \item DFG \citep{bexfield2008}.
\end{itemize}

\textbf{Écarts} (rapportés, non corrigés : la référence Python est normative).
\begin{enumerate}
  \item \ecart{PK du furosémide.} \texttt{fit\_pk.py} fixe la clairance à 0,30\,L/h/kg et vise $t_{1/2}=1{,}25$\,h en l'attribuant à \citet{koh2021}. L'article (lu) donne $735{,}7\pm63{,}8$\,mL/h/kg et $2{,}07\pm0{,}76$\,h. L'ASC IV est $\times2{,}4$.
  \item \ecart{PD du furosémide (placeholder D)} sur 8\,h après 2\,mg/kg :
  \begin{itemize}
    \item diurèse IV 36,8 contre 24,65\,mL/kg ; per os 18,6 contre 30,9 ;
    \item $\Delta$K$^+$ à 8\,h $-1{,}14$ contre $-0{,}33$\,mmol/L ;
    \item $\Delta$Na$^+$ $-8{,}6$ alors que l'article ne relève pas de variation significative ;
    \item $\Delta$Cl$^-$ $-14{,}3$ contre $-4$ à $-6$.
  \end{itemize}
  D'où l'hypokaliémie excessive en prises répétées.
  \item \ecart{Fréquence respiratoire} constante en ICC : 23 contre une médiane de 64/min \citep{yin2026}.
  \item \ecart{VFC en stade C} : SDNN 84 contre $138\pm56$\,ms ; LF/HF 0,25 contre $0{,}40\pm0{,}13$ \citep{bogucki2017}.
  \item \ecart{RSA} : le sens est conforme. La VFC baisse avec la FC, l'atropine et la sévérité \citep{billman2013,baisan2024,keene2019}. L'ampleur n'est pas comparable (bpm contre ms ou \%RSA).
  \item \ecart{Seuil d'œdème} : 18\,mmHg contre 20--24 \citep{yao2025}, citation secondaire de Guyton et Lindsey 1959, non lu.
  \item \ecart{MRC} :
  \begin{itemize}
    \item Na$^+$, Cl$^-$ et K$^+$ constants jusqu'au stade~4 ;
    \item acidose uniquement à trou anionique élevé ;
    \item créatinine de 140\,µmol/L en haut de la plage IRIS1 du module, donc déjà en stade 2 \citep{iris2026}.
  \end{itemize}
  \item \ecart{Ionogramme de la MRC} (recherche dédiée) :
  \begin{itemize}
    \item K$^+$ constant à 4,1 jusqu'au stade~4, alors que les chiens en MRC ont 4,8 contre 4,3 chez les sains, et que 47\,\% font au moins un épisode d'hyperkaliémie \citep{segev2010} ;
    \item HCO$_3^-<22$ seulement à partir du stade~3, contre 43/67/78/76\,\% des chiens aux stades~1 à 4 \citep{lippi2023} ;
    \item phosphate dans l'étendue publiée \citep{parker2017,cortadellas2010} ;
    \item hématocrite trop haute de 4 à 5 points \citep{lippi2021} ;
    \item calcium non modélisé.
  \end{itemize}
  \item \ecart{Hyperkaliémie (Addison), relevé par l'utilisateur} : à K 6,99 la FC est de 106 (le baroréflexe annule le frein \texttt{kfac} de 0,70), alors que la littérature associe hyperkaliémie et bradycardie, avec perte des ondes P et échappements dès K~$\geq$~7. K$^+$ ne modifie que l'aspect de l'ECG, pas la conduction. Correction validée : remplacer \texttt{kfac} par les effets de K$^+$ sur la conduction (proposition rythme).
  \item \ecart{Pimobendane} : T\textsubscript{max} 0,95 contre $2{,}1\pm0{,}9$\,h.
\end{enumerate}

Bases du futur modèle urinaire :
\begin{itemize}
  \item densité urinaire $1{,}040\pm0{,}011$ \citep{rudinsky2019} ;
  \item osmolalité $1541\pm527$\,mOsm/kg \citep{vanvonderen1997} ;
  \item FE Na/Cl/K 0,20/0,42/9,47\,\% \citep{troia2018} ;
  \item DFG \citep{bexfield2008,kelley2022,kelly2016}.
\end{itemize}
Aucune collecte primaire de diurèse sur 24\,h n'a été trouvée.

\section{Provenance des données}
\begin{longtable}{@{}p{3.2cm}p{5.2cm}p{7.6cm}@{}}
\toprule
\textbf{Bloc} & \textbf{Source} & \textbf{Usage et limite}\\
\midrule\endhead
Hémodynamique, gaz & \citet{haskins2005} & FC 87, PAM 103, PAP 14, PAOP 5,5, index cardiaque, gaz artériels et veineux mêlés. Seul le résumé a été lu.\\
Acide-base & \citet{constable2005} & Na, K, Cl, albumine, phosphate ; A\textsubscript{tot} de l'albumine. Coefficient du phosphate humain.\\
Affinité O\textsubscript{2} & \citet{cambier2004} & P50 30,0\,mmHg, effet Bohr $-0{,}37$.\\
VFC & \citet{bogucki2015} & RR moyen, SDNN, RMSSD, LF/HF. Source unique, à recouper.\\
PK furosémide & \citet{koh2021,noah,ema1999} & V\textsubscript{ss}, demi-vie IV, C\textsubscript{max} per os.\\
PK pimobendane & \citet{her2020} & C\textsubscript{max}, ASC, demi-vie. Bras oral 0,5\,mg/kg confirmé (texte intégral lu, $n=7$).\\
MMVD & \citet{keene2019,oyama2004,msdvet} & Stades ACVIM, POG, fraction de régurgitation.\\
MRC & \citet{iris,iris2026} & Bornes de créatininémie par stade IRIS (édition 2026 lue).\\
Addison & \citet{msdvet} & Na:K $<27$.\\
Auscultation & \citet{tiho} & Deux enregistrements (licence CC BY-NC-ND 4.0, aucun embarqué).\\
\bottomrule
\end{longtable}

\section{Statut des références}\label{sec:refs}
Vérification du 17/11/2026 via OpenAlex (titre, auteurs, année, DOI) : toutes les entrées \texttt{@article}, sauf \citet{dormand1980} (vérifié antérieurement). Le champ \texttt{note} de chaque entrée indique jusqu'où la source a été lue (résumé ou texte intégral). \textbf{Documents à préciser} (édition, date de consultation) : \citet{iris,msdvet,noah,ema1999,tiho}. Sources citées dans le rapport mais non vérifiées : Wong 1977 (vue par moteur seulement), Roche-Catholy 2021.

\bibliographystyle{plainnat}
\bibliography{references}
\end{document}
